\documentclass[11pt]{article}

% EMNLP 2026 submission. Uses the official ACL template (acl.sty).
% Switch "review" -> "final" for camera-ready, "preprint" for non-anonymous with page numbers.
\usepackage[review]{acl}

\usepackage{times}
\usepackage{latexsym}
\usepackage[T1]{fontenc}
\usepackage[utf8]{inputenc}
\usepackage{microtype}
\usepackage{inconsolata}
\usepackage{graphicx}

\title{Debias It Yourself: Cognitive Intervention for Bias Mitigation}

\author{
  Anonymous EMNLP submission
}

\begin{document}
\maketitle

\begin{abstract}
Large language models continue to reproduce social bias, and existing mitigations tend to be narrow: a single prompt template, a one-off fine-tune, or a benchmark-specific filter that rarely transfers across tasks, models, or scales. We argue that bias mitigation should instead be grounded in the decades of social-psychology research on \emph{how prejudice is reduced in people}, and operationalized as a reusable behavior the model can be taught and prompted to apply. We introduce \textsc{Debias It Yourself} (DIY), a structured supervision framework that translates established prejudice-reducing interventions into a shared protocol for both training and inference, composing three mechanisms---in-context demonstration, instruction tuning, and intervention-guided self-revision---into a single family of methods. Across multiple bias and reasoning benchmarks and across model families and scales, DIY consistently lowers bias error without degrading task competence, with the largest and most stable gains when supervised teaching is paired with intervention-guided revision. The result is a portable recipe rather than a model-specific trick, and a step toward bias mitigation that is principled, composable, and evaluated as a transferable behavior.
\end{abstract}

\section{Introduction}
Placeholder. Section content will be added in subsequent edits.

% \bibliography{anthology,custom}

\appendix

\section{Appendix}
Placeholder.

\end{document}
